\section{Introduction}
\label{sec:intro}

Behavior-Driven Development (BDD) is an established software engineering methodology that fosters collaboration among software developers, Quality Assurance (QA) engineers, and business stakeholders. By formalizing software requirements into structured, domain-specific natural language scenarios using Gherkin syntax (\texttt{Given-When-Then} clauses), BDD aligns technical test specifications directly with business expectations. Once Gherkin scenarios are defined, automation engineers implement Step Definitions in programming languages such as Python (using Behave) or Java (using Cucumber) to bind step statements to executable web driver actions (e.g., Selenium calls).

Despite its widespread adoption, manually authoring Gherkin scenarios and implementing underlying step definition scripts is labor-intensive, time-consuming, and highly susceptible to human error. As software applications evolve, maintaining comprehensive BDD test suites becomes a severe operational bottleneck.

To address these challenges, researchers have increasingly investigated Large Language Models (LLMs) for automated software test generation. While state-of-the-art foundation models demonstrate strong code generation capabilities, existing studies predominantly restrict evaluation to single-stage Gherkin scenario generation or lack comparative evaluations across different prompt engineering paradigms for multi-artifact BDD generation.

To address this gap, this paper presents an empirical evaluation of \texttt{GPT-4o} (accessed via the cloud-based OpenAI Responses API) for end-to-end BDD test artifact generation. We systematically compare Zero-Shot, Few-Shot, and Chain-of-Thought (CoT) prompting strategies across a benchmark dataset of $N=500$ user stories.

Specifically, this study investigates three Research Questions (RQs):
\begin{itemize}
    \item \textbf{RQ1 (Semantic Similarity)}: Which prompting technique achieves the highest semantic similarity compared with the ground truth?
    \item \textbf{RQ2 (Executable Syntax Validity)}: Which prompting technique produces the highest executable syntax validity?
    \item \textbf{RQ3 (Comparative Effectiveness)}: Are there statistically significant differences among the prompting techniques?
\end{itemize}

This paper makes the following main contributions:
\begin{enumerate}
    \item A large-scale empirical evaluation ($N=500$) measuring semantic similarity, Gherkin syntax validation, Python AST parse rate, and content statistics for joint BDD scenario and step definition generation using GPT-4o.
    \item Rigorous non-parametric statistical hypothesis testing (One-Sample and Pairwise Wilcoxon Signed-Rank tests, Binomial exact tests, and McNemar tests) with effect size quantifications ($d$).
    \item A fully open replication package including prompts, generated raw API outputs, and evaluation scripts to facilitate future research.
\end{enumerate}

The remainder of this paper is organized as follows: Section~\ref{sec:related} discusses related work. Section~\ref{sec:method} details the methodology. Section~\ref{sec:results} presents empirical findings. Section~\ref{sec:discussion} analyzes the results and implications. Section~\ref{sec:threats} addresses threats to validity, and Section~\ref{sec:conclusion} concludes the paper.
